% Imporditud: tools/import_eksperimendid.py
% Allikas: Eesti füüsikaolümpiaadi eksperimendi ülesannete
%   kogu (Rael Kalda, 2023), ülesanne Ü106, lahendus L106.
\setAuthor{EFO žürii}
\setRound{lõppvoor}
\setYear{2012}
\setNumber{P E2}
\setDifficulty{3}
\setTopic{vedelike mehaanika}
\setVahendid{purk veega, tühi joogitops märgistusega, kaks kokkuseotud kumminiiti, niit, mõõtejoonlaud, statiiv}
\setPunktid{14}
\setAllikas{Eksperimendi kogu Ü106, lahendus lk 80}
\setLahendusseis{toores}

\prob{Kumminiit}
Määrake mutriga poldi mass.
\textit{Märkus:} kui joogitopsis veenivoo on märgistuseni, on selles topsis 150 ml vett. Vee tihedus on 1000 kg/m$^3$.

\hint

\solu
Kumminiidi pikenemine on võrdeline kumminiidile mõjuva jõuga. Mõõdame kumminiidi pikenemise õhus $\Delta$$l_1$ = $l_1$ $-$ $l_0$ poldile ja mutrile mõjuva raskusjõu Fo = mg mõjul. Mõõdame kumminiidi pikenemise $\Delta$$l_2$ = $l_2$ $-$ $l_0$, kui polt mutriga on vees Fv = Fo $- \rho$gV . Lahutame teisest seosest esimese. Kumminiidi pikenemine $\Delta$lvesi = $\Delta$$l_1$ $- \Delta$$l_2$ vastab keha poolt välja tõrjutud veele mõjuvale raskusjõule Fo $-$ Fv = $\rho$gV . Arvutame keha poolt välja tõrjutud veele mõjuva raskusjõu mvesi g = $\rho$V g. Sellega oleme kalibreerinud kumminiidi.

Keha ruumala mõõtmiseks tuleb kujundada skaala vee ruumala mõõtmiseks. Me

saame mõõta vee nivoo muutust, me ei saa eriti hästi mõõta purgis oleva veepinna

pindala.

Valame vee topsi kuni jooneni ja mõõdame veenivoo purgis. Nüüd valame topsist

vee purki ja mõõdame uuesti veenivoo purgis. Nivoode erinevus vastab 150 ml vee-

le. Saame arvutada, kui suur vee ruumala muutus purgis vastab vee nivoo muu-

tusele 1 mm võrra.

Mõõdame veenivoo purgis enne keha sukeldamist ja pärast sukeldamist ning ar-

vutame keha ruumala, mis vastab ka välja tõrjutud vee ruumalale.

Seosest $\Delta$lvesi = mvesi g

$\Delta$$l_1$ mkeha g

saame $\Delta$l1mvesi $\Delta$lvesi mkeha = .

Poldi ja mutri mass kokku on 112 g.
\probend
